\documentclass[11pt,letterpaper]{article}
\usepackage[T1]{fontenc}
\usepackage[utf8]{inputenc}
\usepackage{lmodern}
\usepackage[margin=1in,headheight=15pt]{geometry}
\usepackage{amsmath,amssymb,amsthm,mathtools,mathrsfs}
\usepackage{microtype,booktabs,tabularx,array,longtable}
\usepackage{xcolor,fancyhdr,titlesec,needspace,enumitem}
\usepackage[colorlinks=true,linkcolor=blue!35!black,citecolor=green!30!black,urlcolor=blue!45!black]{hyperref}
\usepackage{xurl}
\definecolor{Ink}{HTML}{243747}
\definecolor{Accent}{HTML}{37666D}
\titleformat{\section}{\Large\bfseries\color{Ink}}{\thesection}{.65em}{}
\titleformat{\subsection}{\large\bfseries\color{Accent}}{\thesubsection}{.65em}{}
\pagestyle{fancy}\fancyhf{}
\fancyhead[L]{\small Hilbert geometry at surreal scales}
\fancyhead[R]{\small Splitting, graphs, and least squares}
\fancyfoot[C]{\thepage}
\setlength{\parskip}{.26em}
\setlength{\parindent}{1.1em}
\setlength{\emergencystretch}{3em}
\widowpenalty=10000\clubpenalty=10000
\setlist{itemsep=.15em,topsep=.3em,leftmargin=2em}
\numberwithin{equation}{section}
\newtheorem{theorem}{Theorem}[section]
\newtheorem{proposition}[theorem]{Proposition}
\newtheorem{lemma}[theorem]{Lemma}
\newtheorem{corollary}[theorem]{Corollary}
\theoremstyle{definition}
\newtheorem{definition}[theorem]{Definition}
\newtheorem{example}[theorem]{Example}
\theoremstyle{remark}
\newtheorem{remark}[theorem]{Remark}
\newtheorem{warning}[theorem]{Warning}
\newcommand{\R}{\mathbb R}
\newcommand{\C}{\mathbb C}
\newcommand{\Q}{\mathbb Q}
\newcommand{\Z}{\mathbb Z}
\newcommand{\N}{\mathbb N}
\newcommand{\No}{\mathbf{No}}
\newcommand{\F}{F}
\newcommand{\K}{K}
\newcommand{\HH}[1]{\mathscr H_\Gamma(#1)}
\newcommand{\BB}[2]{\mathscr B_\Gamma(#1,#2)}
\newcommand{\BH}{\mathscr B_\Gamma(H,H)}
\newcommand{\OO}{\mathcal O}
\newcommand{\mm}{\mathfrak m}
\newcommand{\supp}{\operatorname{supp}}
\newcommand{\ran}{\operatorname{ran}}
\newcommand{\res}{\operatorname{res}}
\newcommand{\Graph}{\operatorname{Graph}}
\newcommand{\rank}{\operatorname{rank}}
\newcommand{\ind}{\operatorname{ind}}
\newcommand{\Span}{\operatorname{span}}
\newcommand{\End}{\operatorname{End}}
\newcommand{\GL}{\operatorname{GL}}
\newcommand{\LC}{\operatorname{lc}}
\newcommand{\ip}[2]{\langle #1,#2\rangle}
\newcommand{\norm}[1]{\lVert #1\rVert}
\newcommand{\abs}[1]{\lvert #1\rvert}
\newcommand{\sH}{\mathop{\sum}\nolimits^{\mathrm H}}
\newcommand{\cH}{\mathcal H}
\newcommand{\cG}{\mathcal G}
\newcommand{\Spl}{\operatorname{Spl}}
\newcommand{\Commit}{9a385d3957bdfe3d9ea79f9a524751c90bd2c894}
\newcolumntype{Y}{>{\raggedright\arraybackslash}X}
\hypersetup{pdftitle={Hilbert Geometry at Surreal Scales: Orthogonal Splitting, Amplified Graphs, and Fredholm Least Squares},pdfauthor={Research manuscript prepared for Vladimir Reshetnikov},pdfsubject={Surreal and surcomplex Hahn--Hilbert spaces; orthogonal projections; closed-range threshold; Moore--Penrose inverses}}
\begin{document}
\hypersetup{pageanchor=false}
\begin{titlepage}
\centering
\vspace*{7mm}
{\large\scshape A research manuscript for the Surreal project\par}
\vspace{10mm}
{\Huge\bfseries\color{Ink} Hilbert Geometry\\[.15em] at Surreal Scales\par}
\vspace{7mm}
{\Large Orthogonal Splitting, Amplified Graphs,\\[.2em] and Fredholm Least Squares\par}
\vspace{8mm}
{\large September 23, 2026\par}
\vspace{7mm}
\begin{minipage}{.96\textwidth}
\begin{abstract}\noindent
We develop Hilbert geometry over set-sized surreal and surcomplex Hahn
fields while retaining an ordinary Hilbert space at each coefficient.
The scalar-valued inner product is positive definite, its norm exists
without assuming a divisible value group, and the resulting valued space
is spherically complete. These facts do not imply the projection theorem.

Using the repository's automatic adjoint theorem, we classify every
orthogonally complemented subspace as an infinitesimal unitary deformation
of an ordinary closed subspace. We then prove an exact scale threshold:
for an ordinary bounded operator $T$ and an infinite Hahn scalar $a$,
$\Graph(aT)$ admits orthogonal projection if and only if $\ran T$ is
ordinarily closed; every finite scalar amplification has a complemented
graph. This yields an explicit pair of orthogonally complemented subspaces
whose intersection is not orthogonally complemented. In fact, for a nontrivial
value group, the poset of complemented subspaces is a lattice exactly
when the ordinary coefficient Hilbert space is finite dimensional.

A positive counterpart is an arbitrary-rank Fredholm least-squares
 theorem. Every positive-order Hahn perturbation of an ordinary Fredholm
operator has an adjointable Moore--Penrose inverse. A finite Schur matrix
 determines its kernel, cokernel, index, and exact inverse valuation:
if its rank is $r>0$, then $v(A^\dagger)=-(\nu_r-\nu_{r-1})$, where
$\nu_j$ is the least valuation of its nonzero $j$-minors.
\end{abstract}
\end{minipage}
\vfill
\begin{minipage}{.96\textwidth}\small
\textbf{Scope and status.} The main statements are proved in the text.
The construction, automatic adjoints, direct rotation formula, and
classical perturbation tools are credited, not claimed as new.
The proposed contributions are the exact graph threshold, the
finite-dimensional/lattice dichotomy, and the combined finite-defect
least-squares and sharp-valuation package. A targeted comparison did not
locate these exact statements; priority is not certified. This is an
AI-assisted research draft, not independently refereed or formally
verified. Finite symbolic checks are supplementary, not proofs.

\textbf{Repository pin.} \texttt{9a385d3957bdfe3d9ea79f9a524751c90bd2c894}.
The main arguments are set-sized; Section~\ref{sec:global} separately
treats full surcomplex vector classes. No named published conjecture
is asserted to be solved.
\end{minipage}
\end{titlepage}
\hypersetup{pageanchor=true}
\pagenumbering{roman}
\setcounter{tocdepth}{2}
\tableofcontents
\clearpage
\pagenumbering{arabic}

\section{The question, the model, and the new results}
\label{sec:intro}

Replacing real or complex vector components by surreal or surcomplex
numbers is not merely a change of notation. In infinite dimension one
must choose which coordinate families are vectors, which infinite sums
are legitimate, and which topology defines completeness. Different
choices produce genuinely different theories.

We use the coefficient-Hilbert model already constructed in the Surreal
repository's \emph{Infinite-Dimensional Hahn Spectral Theory}
\cite{RepoSpectral}:
\[
 \cH=\HH{H}=H((t^\Gamma)),\qquad
 \F=\R((t^\Gamma)),\qquad \K=\C((t^\Gamma))=\F[i].
\]
Here $H$ is an ordinary complex Hilbert space and $\Gamma\ne0$ is a
set-sized ordered abelian group. A vector is a well-supported formal
series of ordinary Hilbert vectors. Through a chosen normal-form
embedding of the scalar field, its coordinates are actual surcomplex
numbers. There is also an entirely real version.

This model retains every ordinary vector of $H$, not only vectors with
finite coordinate support. It also retains ordinary bounded operators
as constant coefficients. It does \emph{not} retain ordinary Hilbert
norm convergence as the topology on the constant copy of $H$.
The distinction is crucial both for what works and for what fails.

\subsection{A usable replacement for the projection theorem}
Call a $\K$-linear subspace $M\subseteq\cH$ \emph{orthogonally split}
when $\cH=M\oplus M^\perp$. This is equivalent to every vector having a
nearest point in $M$, and to $M$ being the range of a self-adjoint
idempotent. No prior closedness assumption is required for the
equivalence.

Our first structural theorem says that every such $M$ is exactly
\[
 M=U\HH{M_0},\qquad U^*U=UU^*=I,\qquad v(U-I)>0,
\]
for a uniquely determined ordinary closed $M_0\subseteq H$. There is a
canonical direct-rotation choice of $U$, although there are generally
many possible unitaries. Equivalently, relative to
$H=M_0\oplus M_0^\perp$, the subspace is the graph of a positive-order
Hahn series of bounded operators. This is a classification of all
orthogonally split subspaces, not of all valuation-closed subspaces.

The algebraic rotation formula is classical \cite{Simon}; its use in the
repository for finite spectral clusters is also prior work
\cite{RepoSpectral}. We use it here as a tool for an intrinsic global
subspace classification and the obstructions that follow.

\subsection{An exact graph threshold and a lattice dichotomy}
For an ordinary bounded $T:H\to G$ and a scalar $a\in\K$, let
\[
 M(a,T)=\{(x,aTx):x\in\HH{H}\}
       \subseteq\HH{H\oplus G}.
\]
Theorem~\ref{thm:threshold} proves
\[
 M(a,T)\text{ is orthogonally split}
 \quad\Longleftrightarrow\quad
 \bigl(a=0\text{ or }v(a)\ge0\bigr)
 \ \text{or}\ \ran T\text{ is closed in }G.
\]
For $v(a)<0$, the residue of the graph is precisely
$\ker T\oplus\ran T$. Thus a failure of ordinary closed range becomes a
failure of surreal best approximation after \emph{any} infinite
amplification, however small its negative valuation is.

This has a structural consequence beyond a single counterexample.
Theorem~\ref{thm:lattice} proves that, for every nonzero $\Gamma$,
\[
 \Spl(\HH{H})\text{ is a lattice}
 \quad\Longleftrightarrow\quad \dim_\C H<\infty.
\]
For infinite $H$ it is still an orthomodular poset, with orthogonal sums
and orthocomplementation, but some pairs have neither the relevant meet
nor, after taking complements, the relevant join.

\subsection{A positive theorem: finite defects remain manageable}
Let $T_0:H\to G$ be an ordinary Fredholm operator, and let
$A=T_0+E$ where $E$ is a globally Hahn-supported series of bounded
operators of positive order. A finite matrix $S$ between
$\ker T_0$ and $\ker T_0^*$ controls the entire defect. If
$m=\dim\ker T_0$, $n=\dim\ker T_0^*$, and $r=\rank_\K S$, then
\[
 \dim_\K\ker A=m-r,\qquad
 \dim_\K(\cG/\ran A)=n-r,\qquad
 \ind_\K A=m-n.
\]
Moreover $A^\dagger$ exists in the adjointable Hahn operator algebra,
so all least-squares and minimum-norm problems have exact solutions.
For $r>0$ the largest amplification scale is the explicit invariant
\[
 v(A^\dagger)=-\bigl(\nu_r(S)-\nu_{r-1}(S)\bigr),
 \qquad \nu_0(S)=0.
\]
The definition of $\nu_j$ uses only finitely many minors. No
Archimedean hypothesis, divisibility, analytic convergence, or bound on
all coefficient operator norms is imposed.

\subsection{What this article does not claim}
We do not claim the invention of non-Archimedean inner products,
non-Archimedean failures of Riesz representation or orthogonal
complementation, formal deformation of projections, the Schur
complement, or the Moore--Penrose inverse. Nearby precedents include
Aguayo--Nova's work on $c_0$-type spaces \cite{AN}, Hermitian deformation
theory \cite{BW}, generalized-inverse perturbation theory
\cite{Avrachenkov}, and Hahn-meromorphic Fredholm theory \cite{MS}.
Section~\ref{sec:provenance} distinguishes these from the precise claims
made here.

\section{Scalars, supports, and the surreal interpretation}
\label{sec:scalars}

\subsection{Standing conventions}
An ordered abelian group is written additively. A subset is well ordered
in its given order, not in the reverse order. For a vector space $V$,
\[
 V((t^\Gamma))=
 \left\{x=\sum_{\gamma\in\Gamma}x_\gamma t^\gamma:
 x_\gamma\in V,\ \supp x\text{ well ordered}\right\}.
\]
Set $v(0)=+\infty$ and $v(x)=\min\supp x$ otherwise.
The field $\F$ is ordered by its leading nonzero coefficient, and
complex conjugation on $\K$ acts on coefficients and fixes monomials.
The valuation ring and its maximal ideal are
\[
 \OO=\{a\in\K:v(a)\ge0\},\qquad
 \mm=\{a\in\K:v(a)>0\}.
\]
Zero is included. Residue is extraction of the coefficient of $t^0$.
A scalar is finite in the ordered-field modulus precisely when it is in
$\OO$. An infinite scalar has negative valuation.

\begin{definition}[Strong Hahn summation]
A set-indexed family $(x_j)_{j\in J}$ in $V((t^\Gamma))$ is strongly
Hahn summable if the union of its supports is well ordered and, at each
exponent, only finitely many members have a nonzero coefficient.
The sum $\sH_jx_j$ is formed by these finite coefficient sums.
\end{definition}

We use the classical Hahn--Neumann support lemma: finite sums of
well-ordered subsets of an ordered abelian group are well ordered,
with finitely many decompositions of any fixed exponent. If
$S\subseteq\Gamma_{>0}$ is well ordered, its generated additive monoid
$S^*$ is well ordered, and any exponent is the sum of only finitely many
finite words from $S$. A support-engine proof and attribution are given
in the repository's spectral reports \cite{RepoSpectral}. These facts
justify every formal multiplication and positive-order power series
used below.

In particular, in a unital coefficient algebra,
\begin{equation}\label{eq:neumann}
 (I+B)^{-1}=\sH_{j\ge0}(-B)^j\qquad (v(B)>0).
\end{equation}
Binomial series $(I+B)^q$ for rational $q$ are also defined by the same
support lemma. These are identities of Hahn series. They do not
presuppose topological convergence of the finite partial sums.

\begin{example}[A summable family not converging through finite sums]
In $\C((t^\Q))$, the family $t^{n/(n+1)}$, $n\ge1$, is strongly summable.
Its exponents increase to $1$. The tail of every finite partial sum still
has valuation below $1$, so the partial sums do not converge in the
valuation topology. This distinction is already emphasized in the
repository's foundations and duality reports \cite{RepoFound,RepoDuals}.
\end{example}

\subsection{Actual surreal and surcomplex components}
Fix an order-preserving additive embedding
$\iota:\Gamma\hookrightarrow(\No,+)$. For example, the usual
$\Z,\Q,\R$ can be embedded using their standard surreal copies.
Conway normal forms give the embedding
\begin{equation}\label{eq:embedding}
 \sum_\gamma r_\gamma t^\gamma
 \longmapsto
 \sum_\gamma r_\gamma\omega^{-\iota(\gamma)}.
\end{equation}
The support on the right is reverse well ordered in its exponents,
exactly as a surreal normal form requires. Addition, multiplication,
order, and strong summation are preserved; adjoining $i$ gives an
embedding into $\No[i]$ \cite{Conway,RepoFound}.

For a fixed ordinary orthonormal basis $(e_j)$ of $H$, every vector in
$\HH{H}$ therefore has coordinates in a specified subfield of
$\No[i]$. Section~\ref{sec:coordinates} gives the exact coordinate
admissibility condition. The real construction uses $H$ real and
$\F$ instead of $\K$, with no other change in the main results.

\subsection{Why the whole surreal class is not silently a set field}
A nonzero vector space over all of $\No$ or $\No[i]$ cannot have a
set-sized underlying collection: multiplication of any fixed nonzero
vector by all scalars would inject a proper class into that set.
Thus a literal global theory requires class-sized objects or a
universe-relative formulation. The main theorems below are set-sized
and can be transported by \eqref{eq:embedding}; Section~\ref{sec:global}
then states a separate class-theoretic extension.

For any specified set of scalar normal forms, their supports lie in a
common set-sized additive subgroup of surreal exponents. Thus particular
calculations take place in a common Hahn workspace. This does not mean
that every class-sized problem has been reduced to a single set field,
nor that the full surreal fine topology equals the intrinsic topology of
that workspace.

\section{The Hahn--Hilbert space and its completeness}
\label{sec:construction}

\subsection{The inner product and its norm}
Take ordinary inner products to be linear in the first argument.
For $x,y\in\HH{H}$ put
\begin{equation}\label{eq:inner}
 \ip{x}{y}=\sum_{\alpha,\beta}
           \ip{x_\alpha}{y_\beta}_H t^{\alpha+\beta}.
\end{equation}
At each exponent only finitely many coefficient pairs contribute.
The ordinary Hilbert inner products are evaluated \emph{before} the
Hahn convolution. This is the repository's coefficient-Hilbert inner
product, restated with its proof for the present applications
\cite{RepoSpectral}.

\begin{proposition}[Positive-definite Hahn extension]\label{prop:inner}
The form \eqref{eq:inner} is Hermitian and positive definite. Every
squared norm has a positive square root in $\F$, regardless of whether
$\Gamma$ is divisible. With $\norm{x}=\sqrt{\ip{x}{x}}$,
\[
 v(\norm{x})=v(x),\qquad
 v(\ip{x}{x})=2v(x)\quad(x\ne0).
\]
Cauchy--Schwarz, the triangle inequality, the parallelogram identity, and
finite Pythagoras hold over $\F$.
\end{proposition}
\begin{proof}
Hermitian symmetry and sesquilinearity follow by the finite rearrangement
at each coefficient. If $\delta=v(x)$, the coefficient of
$t^{2\delta}$ in $\ip{x}{x}$ is exactly
$c=\norm{x_\delta}_H^2>0$. Consequently
\[
 \ip{x}{x}=ct^{2\delta}(1+u),\quad v(u)>0\text{ or }u=0,
\]
and its positive root is
\[
 \sqrt c\,t^\delta\sH_{j\ge0}\binom{1/2}{j}u^j.
\]
The support lemma proves existence and the binomial identity proves the
square-root assertion. Positivity makes the root unique. This also
constructs $\abs a=\sqrt{a\bar a}$ for $a\in\K$.
For $y\ne0$ expand the nonnegative scalar
$\ip{x-by}{x-by}$ with $b=\ip{x}{y}/\ip{y}{y}$.
It gives $\abs{\ip{x}{y}}^2\le\ip{x}{x}\ip{y}{y}$.
The triangle inequality follows by expanding $\norm{x+y}^2$ and using
Cauchy--Schwarz. The remaining identities are finite expansions.
\end{proof}

\begin{warning}
The assertion is not that every positive element of $\F$ has a square
root. For example, $t$ has no square root when $\Gamma=\Z$.
Squared norms have even leading valuation, which is exactly why all roots
needed here do exist.
\end{warning}

\subsection{Valuation topology and spherical completeness}
A neighborhood basis of $0$ is $\{x:v(x)>\gamma\}$, $\gamma\in\Gamma$.
It is the same topology as that generated by all positive
$\F$-valued norm radii: a strict increase of valuation beats every
ordinary leading coefficient, and the positive monomial radii are
cofinal among the needed tests. In rank one, when $\Gamma\subseteq\R$,
$d(x,y)=\exp(-v(x-y))$, with $d(x,x)=0$, is a real-valued ultrametric
for the same topology. It is not the $\F$-valued Hilbert norm itself.

\begin{theorem}[Spherical completeness of the coefficient space]
\label{thm:complete}
For every ordinary vector space $V$, every nonempty nested set of balls
\[
 B(a,\gamma)=\{x:v(x-a)\ge\gamma\}
       \quad(a\in V((t^\Gamma)),\ \gamma\in\Gamma)
\]
has nonempty intersection. In particular $V((t^\Gamma))$ is complete
for valuation-Cauchy nets. Thus $\HH{H}$ is spherically complete even
when its orthogonal projection theorem fails.
\end{theorem}
\begin{proof}
A ball fixes all coefficients strictly below its radius exponent.
Nested balls impose compatible fixed coefficients. At an exponent fixed
by at least one ball, choose that forced value; at all other exponents
choose zero. Call the resulting coefficient family $x$.

Its support is well ordered. Indeed, if it contained a strictly
decreasing sequence $\delta_1>\delta_2>\cdots$, some ball fixes the
nonzero coefficient at $\delta_1$ and hence every coefficient at
$\delta_j$. The entire sequence would then occur in the well-ordered
support of that ball's center, a contradiction. In ZFC the absence of
such a sequence in a linearly ordered set is equivalent to
well-ordering. Thus $x$ is a Hahn vector and belongs to every ball.

For completeness, a valuation-Cauchy net has an eventually fixed
truncation below every exponent. These truncations are compatible.
The same argument gives a Hahn vector with those truncations, and the
net converges to it. The argument uses no norm convergence in $V$.
\end{proof}

The last sentence is informative: completeness here does not force
ordinary completeness on coefficient subspaces. If $V\subset H$ is
any ordinary linear subspace, even one that is not ordinarily closed,
then $V((t^\Gamma))$ is valuation closed in $H((t^\Gamma))$. To see this,
a coefficient outside $V$ cannot be changed by a sufficiently small
valuation perturbation. This is one source of the projection failure.

\subsection{Finite-dimensional Hilbert geometry survives}
\begin{proposition}\label{prop:finite}
Every finite-dimensional $\K$-linear subspace of $\HH{H}$ has an
orthonormal basis and an orthogonal projection. If
$u_1,\ldots,u_d$ is such a basis, then
\[
 P_Mx=\sum_{j=1}^d\ip{x}{u_j}u_j.
\]
The coefficient-zero vectors $(u_j)_0$ form an ordinary orthonormal
family, and the residue of $M$ has ordinary dimension $d$.
\end{proposition}
\begin{proof}
Finite Gram--Schmidt uses only positive squared norms and their roots
from Proposition~\ref{prop:inner}. The displayed formula is then a
self-adjoint idempotent onto $M$. Since $\norm{u_j}=1$, one has
$v(u_j)=0$, and taking constant coefficients of the orthogonality
identities gives ordinary orthonormality.
For $x=\sum a_ju_j$, positivity gives
$v(x)=\min_jv(a_j)$; there is no cancellation of its least coefficient
because the leading coefficient vectors of the $u_j$ are orthonormal.
Thus integral vectors in $M$ have precisely the residue span of those
$d$ vectors.
\end{proof}

\section{Coordinates, two summation levels, and duality}
\label{sec:coordinates}

\subsection{An exact coordinate-space description}
Choose an ordinary orthonormal basis indexed by a set $I$, identifying
$H$ with $\ell^2(I,\C)$. The coordinate map sends
$x=\sum x_\gamma t^\gamma$ to
$a_i=\sum_\gamma(x_\gamma)_i t^\gamma$.

\begin{proposition}[Admissible surcomplex coordinate arrays]
The coordinate map identifies $\HH{\ell^2(I)}$ with the families
$(a_i)\in\K^I$ satisfying both conditions:
\[
 \bigcup_{i\in I}\supp a_i\text{ is well ordered},\qquad
 \bigl([t^\gamma]a_i\bigr)_{i\in I}\in\ell^2(I)
       \quad\text{for every }\gamma.
\]
The inner product is
\begin{equation}\label{eq:twosums}
 [t^\delta]\ip{x}{y}
 =\sum_{\alpha+\beta=\delta}
     \left(\sum_{i\in I}(x_\alpha)_i
                         \overline{(y_\beta)_i}\right).
\end{equation}
The outer sum is finite and each inner sum is ordinarily absolutely
convergent.
\end{proposition}
\begin{proof}
The conditions are precisely that the coefficient rows define ordinary
Hilbert vectors and that their nonzero support is a Hahn support.
They provide the inverse coordinate map. Cauchy--Schwarz in the ordinary
$\ell^2(I)$ proves absolute convergence of the inner sum; the support
lemma proves finiteness of the outer sum.
\end{proof}

\begin{example}[Why a single formal coordinate sum is inadequate]
The constant vector $x=(1/n)_{n\ge1}$ belongs to $\HH{\ell^2}$ and has
its usual real squared norm. But the scalar family $(1/n^2)_{n\ge1}$
is not strongly Hahn summable: infinitely many nonzero terms occur at
exponent zero. Formula~\eqref{eq:twosums} uses ordinary Hilbert summation
at a fixed coefficient, not that inadmissible formal sum. Conversely,
$x=\sum_{n\ge1}e_n t^{n/(n+1)}$ is a legitimate Hahn vector even though
its monomial partial sums do not converge in valuation.
\end{example}

Thus ``square summable surcomplex coordinates'' must include a summation
rule. Merely requiring finite partial squared sums to be bounded is
far too weak: an infinite scalar already bounds every ordinary integer,
so the constant family of coordinate values $1$ would pass that test.

\subsection{What becomes of an orthonormal basis}
The constant $e_i$ remain orthonormal and separate all vectors by inner
products. They need not be a valuation-topological Schauder basis.
A nonzero ordinary tail, however small its ordinary Hilbert norm, has
valuation zero. Its size therefore does not tend to zero at all surreal
scales. Basis synthesis in this model is \emph{layered}: ordinary
Hilbert synthesis within each coefficient, then Hahn assembly across
exponents. Parseval holds coefficientwise in the sense of
\eqref{eq:twosums}, not as an unrestricted strong sum over $i$.

\subsection{Riesz representation requires more than valuation continuity}
The repository's three-duals report \cite{RepoDuals} proves a detailed
classification. One elementary obstruction is enough here. In infinite
dimension choose an everywhere-defined but ordinarily discontinuous
complex-linear functional $\ell:H\to\C$ using a Hamel basis, and set
\[
 L\left(\sum x_\gamma t^\gamma\right)
       =\sum\ell(x_\gamma)t^\gamma.
\]
It respects strong sums and satisfies $v(Lx)\ge v(x)$, hence is
valuation continuous. It cannot equal $\ip{x}{y}$ for a Hahn vector
$y$: testing constants forces all nonconstant coefficients of $y$ to
vanish and represents $\ell$ by an ordinary Hilbert vector, a
contradiction.

For a fixed $\eta>0$, this functional even satisfies
$\abs{Lx}\le t^{-\eta}\norm{x}$. The infinite bound dominates each
individual ordinary coefficient ratio. Therefore the usual phrase
``bounded functional'' does not restore Riesz unless the meaning of
boundedness is strengthened. By contrast, a globally Hahn-supported
series of ordinary continuous functionals is represented, coefficient
by coefficient, by a unique Hahn vector. Neither observation is claimed
as new here.

\section{The intrinsic algebra of adjointable operators}
\label{sec:adjoints}

\subsection{Operators with a common Hahn support}
For ordinary Hilbert spaces $H,G$, put
\[
 \BB{H}{G}=\mathcal B(H,G)((t^\Gamma)).
\]
An element $A=\sum A_\gamma t^\gamma$ has one well-ordered support of
nonzero ordinary bounded operators. It acts on all Hahn vectors by
\[
 (Ax)_\delta=\sum_{\gamma+\alpha=\delta}A_\gamma x_\alpha.
\]
The support lemma makes these coefficient sums finite. Composition is
Hahn convolution and
$A^*=\sum A_\gamma^*t^\gamma$. There is no requirement that the numbers
$\norm{A_\gamma}$ be uniformly bounded as $\gamma$ varies.
Write $v(A)=\min\supp A$ for $A\ne0$, and $v(0)=+\infty$.
Then
\begin{equation}\label{eq:opvaluation}
 v(Ax)\ge v(A)+v(x),\qquad v(AB)\ge v(A)+v(B).
\end{equation}
An operator is called \emph{integral} if its valuation is nonnegative.
An integral operator with an integral two-sided inverse preserves vector
valuations exactly. Multiplication on either side by such an operator
also preserves the valuation of any nonzero operator: apply
\eqref{eq:opvaluation} first to the operator and then to its inverse.

The apparent restriction to globally supported series is intrinsic to
adjoints, rather than an optional technical convenience. We give the
repository's theorem and proof in the rectangular form needed below.

\begin{theorem}[Automatic adjoint structure; repository input]
\label{thm:adjoint}
Let $A:\HH{H}\to\HH{G}$ and $B:\HH{G}\to\HH{H}$ be everywhere-defined
$\K$-linear maps satisfying
\[
 \ip{Ax}{y}=\ip{x}{By}\qquad(x\in\HH{H},\ y\in\HH{G}).
\]
There is a unique $A\in\BB{H}{G}$ producing the given map, and
$B=A^*$ coefficientwise. In particular every everywhere-defined
adjointable map is strong, valuation continuous, and order bounded.
\end{theorem}
\begin{proof}
The square version and its argument are
\texttt{ihs:hh:thm:adjoint} in \cite{RepoSpectral}. The same argument
works between different coefficient spaces, as follows.

For constants $u\in H,w\in G$, define
$A_\gamma u=[t^\gamma]Au$ and $B_\gamma w=[t^\gamma]Bw$.
Extracting coefficients gives
\[
 \ip{A_\gamma u}{w}_G=\ip{u}{B_\gamma w}_H.
\]
Each $A_\gamma$ has a closed ordinary graph. Indeed, if
$u_j\to u$ and $A_\gamma u_j\to z$ in ordinary Hilbert norms, testing
against an arbitrary $w$ gives
$\ip{z}{w}=\ip{u}{B_\gamma w}=\ip{A_\gamma u}{w}$.
Thus $z=A_\gamma u$. The ordinary closed graph theorem makes
$A_\gamma$ bounded. The same argument applies to $B_\gamma$, and the
identity identifies it with $A_\gamma^*$.

Suppose the set $S=\{\gamma:A_\gamma\ne0\}$ were not well ordered.
Choice gives a strictly decreasing sequence
$\gamma_1>\gamma_2>\cdots$ in $S$. Each $\ker A_{\gamma_j}$ is a proper
closed subspace of $H$, hence nowhere dense. By the ordinary Baire
category theorem some $u$ lies outside their union. All $\gamma_j$
would then lie in the support of the Hahn vector $Au$, a contradiction.
Therefore $S$ is well ordered. The adjoint coefficients have the same
support.

Form the convolution operators $\widetilde A,\widetilde B$ from these
coefficient families. They are adjoints, and agree with $A,B$ on
constants. For arbitrary $x$ and constant $w$,
\[
 \ip{Ax}{w}=\ip{x}{Bw}
 =\ip{x}{\widetilde B w}=\ip{\widetilde A x}{w}.
\]
Pairing with every constant vector separates Hahn coefficients, so
$Ax=\widetilde A x$. The analogous argument proves
$B=\widetilde B$. This also proves uniqueness.

The support lemma shows that convolution preserves strong sums, and
\eqref{eq:opvaluation} gives continuity. For order boundedness let
$s=v(A)$ and choose an ordinary real $c>\norm{A_s}$. For nonzero
$x$ with leading term $t^\delta u$, the coefficient at
$2(s+\delta)$ of
$c^2t^{2s}\norm{x}^2-\norm{Ax}^2$ is
$c^2\norm{u}^2-\norm{A_su}^2>0$. Hence
$\norm{Ax}\le ct^s\norm{x}$.
\end{proof}

\begin{remark}
The theorem uses ordinary Hilbert completeness in its coefficient
closed-graph and Baire arguments. Spherical completeness of the Hahn
space alone does not supply those arguments for arbitrary coefficient
vector spaces. Nor does the theorem classify all valuation-continuous
maps: the three-duals report gives strictly larger classes
\cite{RepoDuals}.
\end{remark}

\subsection{Finite-rank operators and self-adjoint inverses}
For fixed Hahn vectors $u,w$, the map $x\mapsto\ip{x}{w}u$ is
adjointable, with adjoint $y\mapsto\ip{y}{u}w$. Expanding its two
vector supports gives a Hahn series of bounded finite-rank coefficient
operators. Finite sums of these maps include every finite-dimensional
orthogonal projection constructed in Proposition~\ref{prop:finite}.

If a self-adjoint everywhere-defined map $B$ is bijective, then its
algebraic inverse is self-adjoint: writing $x=Bu,y=Bv$ gives
\[
 \ip{B^{-1}x}{y}=\ip{u}{Bv}=\ip{Bu}{v}
 =\ip{x}{B^{-1}y}.
\]
Theorem~\ref{thm:adjoint} then puts $B^{-1}$ in the Hahn operator
algebra. We shall use this specific observation, rather than assume
without proof a general bounded-inverse theorem over $\K$.

\section{Orthogonal splitting and its residue normal form}
\label{sec:splitting}

\subsection{Three equivalent formulations}

\begin{definition}
A $\K$-linear subspace $M\subseteq\cH$ is \emph{orthogonally split}
if $\cH=M\oplus M^\perp$, where
$M^\perp=\{x:\ip{x}{m}=0\text{ for every }m\in M\}$.
The word ``split'' below always means this orthogonal property.
\end{definition}

\begin{proposition}[Projection and best approximation]
\label{prop:projection}
For a subspace $M\subseteq\cH$, the following are equivalent:
\begin{enumerate}[label=\textup{(\roman*)}]
\item $M$ is orthogonally split;
\item $M$ is the range of an everywhere-defined self-adjoint idempotent;
\item every $x\in\cH$ has a nearest point in $M$, with distances
compared in the ordered field $\F$.
\end{enumerate}
The projection and all nearest points are unique. A split subspace is
valuation closed and closed under all strong sums of its vectors.
\end{proposition}
\begin{proof}
A decomposition $x=m+n$ with $m\in M,n\in M^\perp$ is unique by
positivity. Sending $x$ to $m$ defines a linear self-adjoint idempotent.
Conversely the range and kernel of a self-adjoint idempotent are
orthogonal, and its kernel equals the orthogonal complement of its
range.

For such a decomposition and any $m'\in M$, Pythagoras gives
\[
 \norm{x-m'}^2=\norm{n}^2+\norm{m-m'}^2.
\]
Thus $m$ is the unique nearest point. Conversely, let $m$ be nearest
to $x$ and put $r=x-m$. If $n\in M$ and $\ip{r}{n}\ne0$, then $n\ne0$
and the scalar $a=\ip{r}{n}/\ip{n}{n}$ gives
\[
 \norm{r-an}^2=\norm{r}^2-
          \frac{\abs{\ip{r}{n}}^2}{\norm{n}^2}<\norm{r}^2,
\]
contradicting minimality. Hence $r\in M^\perp$ and the splitting
follows. Finally $M=\ker(I-P)$ for its projection $P$; automatic
adjoint structure makes $P$ continuous and strong. These two
properties prove the two closure assertions.
\end{proof}

\subsection{Every projection has ordinary residue}

\begin{lemma}[Integral projections]\label{lem:integralproj}
Every self-adjoint idempotent $P$ in $\BH$ is integral. If $P\ne0$,
then $v(P)=0$. Its residue $P_0$ is an ordinary orthogonal projection,
and $P-P_0$ has strictly positive order or is zero.
\end{lemma}
\begin{proof}
Let $s=v(P)$ when $P\ne0$. Its nonzero leading coefficient $P_s$ is
ordinary self-adjoint, so $P_s^2\ne0$: otherwise
$\norm{P_su}^2=\ip{P_s^2u}{u}=0$ for all $u$.
Consequently $v(P^2)=2s$. The identity $P^2=P$ forces $2s=s$, hence
$s=0$. Taking residues gives $P_0^2=P_0=P_0^*$.
\end{proof}

For any subspace define its \emph{integral residue space} by
\[
 \res M=\{x_0:x\in M,\ v(x)\ge0\}\subseteq H.
\]
This is an ordinary complex-linear subspace, without a presumed
ordinary closedness property.

\begin{corollary}[Closed-residue obstruction]\label{cor:residue}
If $M$ is split with projection $P$, then
\[
 \res M=\ran P_0.
\]
In particular its residue space is ordinarily closed in $H$.
\end{corollary}
\begin{proof}
For integral $x=Px$ one has $x_0=P_0x_0$. Conversely, for constant
$u\in H$, the vector $Pu$ is integral, belongs to $M$, and has residue
$P_0u$. The range of an ordinary orthogonal projection is closed.
\end{proof}

Ordinary closedness in this corollary is a substantially stronger
restriction than valuation closedness of $M$. It will detect the
amplified-graph obstruction without an appeal to spectral theory.

\subsection{A global normal form using the classical direct rotation}

\begin{theorem}[Residue normal form]\label{thm:normalform}
A subspace $M\subseteq\HH{H}$ is orthogonally split if and only if
there are an ordinary closed subspace $M_0\subseteq H$ and a Hahn
unitary $U\in\BH$ satisfying $v(U-I)>0$ or $U=I$, such that
\begin{equation}\label{eq:normalform}
 M=U\HH{M_0}.
\end{equation}
Here $M_0=\res M$ is unique. Given the projection $P$ onto $M$ and
$P_0$ onto $M_0$, a canonical choice is
\begin{align}
 W&=PP_0+(I-P)(I-P_0),\label{eq:rotationW}\\
 D&=I-(P-P_0)^2,\nonumber\\
 U&=W D^{-1/2}.\label{eq:rotation}
\end{align}
The inverse square root is the formal binomial series about $I$.
\end{theorem}
\begin{proof}
For a split $M$, Lemma~\ref{lem:integralproj} gives
$P=P_0+\text{positive-order terms}$. Direct multiplication, using
$P^2=P$ and $P_0^2=P_0$, gives
\[
 W^*W=WW^*=D,\qquad WP_0=PW.
\]
The element $(P-P_0)^2$ commutes with both projections, and thus with
$W$ and $W^*$. Its support is positive, unless it vanishes.
The Hahn--Neumann lemma therefore defines the binomial series
$D^{-1/2}$ and verifies its identities algebraically. It commutes
with the same operators. Consequently
\[
 U^*U=UU^*=I,\qquad UP_0=PU,
\]
and $U=I+\text{positive-order terms}$. It follows that
$P=UP_0U^*$ and hence \eqref{eq:normalform} holds with
$M_0=\ran P_0$. Corollary~\ref{cor:residue} proves uniqueness of
$M_0$. Conversely, the operator $UP_0U^*$ is an orthogonal projection
onto the subspace in \eqref{eq:normalform}.

The formula \eqref{eq:rotation} is the classical direct rotation of
projections; see Simon's account and attribution to earlier work
\cite{Simon}. The point here is that automatic adjoints and integral
projections place \emph{every} split Hahn subspace in its domain of
formal applicability. The formula itself is not a new result.
\end{proof}

\subsection{Exact graph coordinates for a fixed residue}

\begin{corollary}[The positive-order graph chart]\label{cor:graphchart}
Fix an ordinary closed $M_0\subseteq H$ and decompose
$H=M_0\oplus M_0^\perp$. The split subspaces with residue $M_0$ are in
bijection with
\[
 X\in\BB{M_0}{M_0^\perp},\qquad X=0\text{ or }v(X)>0,
\]
by $X\mapsto\Graph X$. For $J=(I,X)^{\mathsf T}$ the projection is
\begin{equation}\label{eq:graphproj}
 P_X=J(I+X^*X)^{-1}J^*
 =\begin{pmatrix}R&RX^*\\XR&XRX^*\end{pmatrix},
 \qquad R=(I+X^*X)^{-1}.
\end{equation}
\end{corollary}
\begin{proof}
Take the near-identity unitary $U$ in Theorem~\ref{thm:normalform}
and write its blocks relative to the displayed decomposition.
The block $U_{00}=I+\text{positive-order terms}$ is invertible by
\eqref{eq:neumann}. Its image of $M_0$ is therefore the graph of
$X=U_{10}U_{00}^{-1}$, a positive-order operator. A graph determines
its map uniquely.

Conversely $I+X^*X$ has a Neumann inverse of residue $I$. Formula
\eqref{eq:graphproj} is self-adjoint and idempotent, since
$J^*J=I+X^*X$; its range is exactly $\Graph X$. Its residue is the
projection onto the first summand. All products and inverses have
support in the positive monoid generated by the supports of $X,X^*$.
\end{proof}

\begin{remark}[No assertion that closed residue is sufficient]
The closed-residue test in Corollary~\ref{cor:residue} is necessary,
not a general characterization of arbitrary closed subspaces. The
complete characterization above also requires an adjointable
positive-order graph chart. We make no claim that an arbitrary
valuation-closed subspace with closed residue automatically has one.
\end{remark}

\section{Amplified graphs: an exact closed-range threshold}
\label{sec:graphs}

\subsection{The normal equation is the whole obstruction}

\begin{proposition}[General adjointable graph criterion]
\label{prop:graphcriterion}
Let $A\in\BB{H}{G}$. Its graph is valuation closed and closed under
strong Hahn sums. It is orthogonally split if and only if
$I+A^*A$ is surjective on $\HH{H}$. In that case this operator is
invertible in $\BB{H}{H}$ and the graph projection is
\[
 P=\begin{pmatrix}I\\A\end{pmatrix}
       (I+A^*A)^{-1}\begin{pmatrix}I&A^*\end{pmatrix}.
\]
\end{proposition}
\begin{proof}
The graph is the kernel of the continuous strong map
$(x,y)\mapsto y-Ax$, which gives both closure properties.
The operator $B=I+A^*A$ is injective because
\[
 \ip{Bx}{x}=\norm{x}^2+\norm{Ax}^2>0\quad(x\ne0).
\]
If the graph is split, project $(f,0)$ onto a point $(x,Ax)$.
Orthogonality of the residual to every $(h,Ah)$ gives $Bx=f$.
As $f$ was arbitrary, $B$ is onto. Conversely, if $B$ is onto it is
bijective and self-adjoint, so the observation following
Theorem~\ref{thm:adjoint} gives $B^{-1}$ in the adjointable algebra.
Writing $Jx=(x,Ax)$, the identity $J^*J=B$ verifies that
$JB^{-1}J^*$ is an orthogonal projection with range $\Graph A$.
\end{proof}

Positivity of $I+A^*A$ does not alone imply its surjectivity. This is
exactly the step in an ordinary Hilbert-space graph argument that
cannot be imported without further hypotheses.

\subsection{Finite and infinite scalar amplification}

\begin{theorem}[Exact amplification threshold]\label{thm:threshold}
Let $T:H\to G$ be an ordinary bounded operator and $a\in\K$.
Then
\begin{equation}\label{eq:threshold}
 \Graph(aT)\text{ is split}
 \quad\Longleftrightarrow\quad
 a=0\text{ or }v(a)\ge0\text{ or }\ran T\text{ is ordinarily closed}.
\end{equation}
For nonzero infinite $a$, one has the exact residue formula
\begin{equation}\label{eq:graphresidue}
 \res\Graph(aT)=\ker T\oplus\ran T\subseteq H\oplus G.
\end{equation}
When the graph is split, its projection has residue the ordinary
projection onto $\Graph(a_0T)$ in the finite case, and onto
$\ker T\oplus\ran T$ in the infinite case.
\end{theorem}
\begin{proof}
If $a=0$ the graph is horizontal. Suppose $v(a)\ge0$, and put
$a_0=[t^0]a$. The constant coefficient of
$B=I+\bar a a T^*T$ is
$B_0=I+\abs{a_0}^2T^*T$, an ordinary boundedly invertible positive
operator. Factor
$B=B_0(I+B_0^{-1}(B-B_0))$ and apply the positive-order Neumann
series. Thus $B$ is invertible in the Hahn operator algebra, and
Proposition~\ref{prop:graphcriterion} proves splitting. Taking the
residue of its projection formula gives the finite-case assertion.

Now let $v(a)<0$. Set $b=a^{-1}$ and $\eta=v(b)>0$, with
$b_\eta\ne0$. A graph vector $(x,y)$ is integral exactly when both
components are integral. It satisfies $Tx=by$. The constant
coefficient gives $Tx_0=0$. The coefficient at $\eta$ gives
\[
 T x_\eta=b_\eta y_0,
\]
because $y$ has no negative exponents and $b$ has least exponent
$\eta$. Hence $(x_0,y_0)\in\ker T\oplus\ran T$.
Conversely, for any $k\in\ker T,u\in H$,
\[
 x=k+bu,\qquad y=Tu
\]
is an integral graph vector with residue $(k,Tu)$. This proves
\eqref{eq:graphresidue}. If the graph is split,
Corollary~\ref{cor:residue} makes this residue space ordinarily
closed, forcing $\ran T$ to be closed in $G$.

For the converse assume $\ran T$ closed. Write
\[
 H=N\oplus X,\quad N=\ker T,\qquad
 G=Y\oplus C,\quad Y=\ran T.
\]
The ordinary operator $T_1=T|_X:X\to Y$ is a bounded isomorphism,
by the ordinary bounded inverse theorem. Let $C_1=T_1^*T_1$ on $X$;
it is boundedly invertible. On $N$, $B=I+\bar a a T^*T$ equals $I$.
On $X$ an inverse is
\begin{equation}\label{eq:infinitegraphinverse}
 B_X^{-1}=(\bar a a)^{-1}C_1^{-1}
      \bigl(I+(\bar a a)^{-1}C_1^{-1}\bigr)^{-1}.
\end{equation}
The inner correction has positive valuation $2\eta$, so its inverse
is Hahn defined. This supplies a two-sided inverse to $B$ and proves
splitting. The residue projection is the ordinary projection onto
\eqref{eq:graphresidue}, by Corollary~\ref{cor:residue}.
\end{proof}

\begin{remark}[Geometric interpretation]
At a finite nonzero residue $a_0$, the leading geometry is the ordinary
sloped graph. If $v(a)>0$, it is horizontal. At an infinite scale and
with closed range, the nonkernel portion of the graph becomes vertical:
its residue is $\ker T\oplus\ran T$. A nonclosed ordinary range
cannot be the residue of an orthogonal projection, so precisely this
case obstructs the transition. The threshold depends on the sign of
$v(a)$, not on a metric magnitude or a discrete value-group step.
\end{remark}

\subsection{An explicit complete graph with no nearest point}

\begin{example}\label{ex:badgraph}
Let $H=G=\ell^2(\N_{\ge1})$, let $Te_n=n^{-1}e_n$, choose
$\eta\in\Gamma_{>0}$, and put $q=t^\eta$, $A=q^{-1}T$.
The range of $T$ is not ordinarily closed: it contains all
finite-support vectors but does not contain $h=(1/n)_{n\ge1}$.
Indeed $Tx=h$ would require $x=(1,1,\ldots)\notin\ell^2$.
The graph of $A$ is therefore not split.

There is also a direct witness to failure of best approximation.
The vector $z=(h,0)$ has no nearest point in the graph. Such a point
$(x,Ax)$ would solve the normal equation
\begin{equation}\label{eq:badnormal}
 (q^2I+T^2)x=q^2h.
\end{equation}
Because $T^2$ is injective, the left side has valuation $v(x)$ for
nonzero $x$. Thus $v(x)=2\eta$. Its coefficient there would satisfy
$T^2x_{2\eta}=h$, requiring
$x_{2\eta}=(n)_{n\ge1}$, again not a Hilbert vector.
The contradiction proves the assertion without a compact-operator
spectral theorem.

This graph is valuation closed, strong-sum closed, and complete as a
closed subspace of the complete valued space in Theorem~\ref{thm:complete}.
These properties do not supply a nearest point. We assert failure of
an attained minimum here; no assertion about a distance infimum is
needed for the example.
\end{example}

\begin{corollary}[Enlarging the scale group cannot repair this graph]
\label{cor:norepair}
Let $\Gamma\hookrightarrow\Delta$ be an ordered group embedding,
and extend coefficients and exponents to the $\Delta$ Hahn spaces.
For ordinary $T$ and $v(a)<0$, the extended graph is split exactly
when the original one is split, namely when $\ran T$ is ordinarily
closed. This holds also for noncofinal embeddings of value groups.
\end{corollary}
\begin{proof}
The image of $a$ still has negative valuation. Theorem~\ref{thm:threshold}
applies over both groups with the same ordinary range condition.
\end{proof}

\section{The orthogonal-subspace poset is usually not a lattice}
\label{sec:lattice}

\subsection{What survives: orthocomplements and orthogonal sums}
Let $\Spl(\cH)$ be the collection of split $\K$-linear subspaces,
ordered by inclusion. Its bottom and top are $0$ and $\cH$.
Orthocomplementation is an order-reversing involution on this
collection, by the defining orthogonal decompositions.

\begin{proposition}[Orthomodular poset structure]\label{prop:poset}
If $M,N\in\Spl(\cH)$ and $M\perp N$, then $M\oplus N$ is split
and is their least upper bound in $\Spl(\cH)$. If $M\subseteq N$,
then $N\cap M^\perp$ is split and
\[
 N=M\oplus(N\cap M^\perp).
\]
Consequently $\Spl(\cH)$ is an orthomodular poset.
\end{proposition}
\begin{proof}
Let $P_M,P_N$ be the unique projections. Orthogonality gives
$P_MP_N=P_NP_M=0$, so $P_M+P_N$ is a self-adjoint idempotent with
range $M\oplus N$. Every subspace containing both contains their sum,
which proves the least-upper-bound assertion.

If $M\subseteq N$, then $P_NP_M=P_M$ and, taking adjoints,
$P_MP_N=P_M$. The operator $P_N-P_M$ is a self-adjoint idempotent.
Its range is exactly $N\cap M^\perp$, proving both splitting and the
displayed orthomodular identity. These properties, together with
orthocomplementation and the bounds $0,\cH$, are the asserted
orthomodular-poset axioms. They do not assert that arbitrary binary
meets or joins exist.
\end{proof}

\subsection{Two good subspaces with a bad intersection}

\begin{theorem}[Exact finite-dimensional/lattice dichotomy]
\label{thm:lattice}
Let $\Gamma\ne0$ be any set-sized ordered abelian group. For an
ordinary real or complex Hilbert space $H$,
\[
 \Spl(\HH{H})\text{ is a lattice}
 \quad\Longleftrightarrow\quad \dim H<\infty.
\]
If $H$ is infinite dimensional, there are two split subspaces whose
intersection is not split and which have no greatest lower bound in
$\Spl(\HH{H})$. Their orthogonal complements have no least upper bound.
\end{theorem}
\begin{proof}
First work in $\HH{H_1\oplus H_1\oplus H_1}$ with
$H_1=\ell^2(\N_{\ge1})$. Use $T,q$ from Example~\ref{ex:badgraph}
and define
\begin{align*}
 M&=\{(x,y,0):x,y\in\HH{H_1}\},\\
 N&=\{(x,y,Tx-qy):x,y\in\HH{H_1}\}.
\end{align*}
The first is a constant orthogonal summand. The second is the graph
of the integral adjointable operator
$B(x,y)=Tx-qy$. The residue of $I+B^*B$ is the ordinary operator
$I+B_0^*B_0$, boundedly invertible because it is at least $I$.
A positive-order Neumann inverse and
Proposition~\ref{prop:graphcriterion} show that $N$ is split.

Their intersection is
\begin{equation}\label{eq:intersection}
 M\cap N=\{(x,q^{-1}Tx,0):x\in\HH{H_1}\}.
\end{equation}
Its residue space is $0\oplus\ran T\oplus0$, since $\ker T=0$.
That space is not ordinarily closed. Corollary~\ref{cor:residue}
therefore proves that the intersection is not split.

Suppose a greatest lower bound $L$ of $M,N$ existed in the split
poset. For every nonzero $z\in M\cap N$, the one-dimensional space
$\K z$ is split by finite-dimensional orthogonal projection, and is
a common lower bound. Greatestness implies $\K z\subseteq L$ for
all such $z$. It follows that $M\cap N\subseteq L$, whereas the
lower-bound condition gives $L\subseteq M\cap N$. This would make
the nonsplit intersection itself split, a contradiction. An
order-reversing involution exchanges meets and joins, so
$M^\perp,N^\perp$ have no least upper bound.

Every infinite-dimensional Hilbert space contains a closed isometric
copy of $H_1\oplus H_1\oplus H_1$, obtained by partitioning a countable
orthonormal family into three infinite families and taking their
closed spans. Extend the two projections by zero on the ordinary
orthogonal complement. They give split subspaces in $\HH{H}$ with
the same intersection. Its residue is still nonclosed in $H$, and
the preceding no-meet argument applies unchanged. All choices here
are set-sized.

Finally, if $\dim H=d<\infty$, then $\HH{H}$ is isometric to
$\K^d$ (or $\F^d$ in the real case). Every linear subspace has a
finite orthonormal basis by Proposition~\ref{prop:finite}, hence is
split. Ordinary algebraic intersection and sum are meet and join
in the full subspace lattice, so the split poset is a lattice.
\end{proof}

\begin{remark}[Relation to existing non-Archimedean obstructions]
The general phenomenon of nonorthomodularity in other non-Archimedean
Hilbert-like models is not new; Aguayo--Nova give such results for
$c_0$-type spaces \cite{AN}. Here the coefficient-Hilbert model,
the complete classification of split subspaces, and the precise
criterion ``finite ordinary coefficient dimension'' identify a
particular obstruction. In particular, the poset of \emph{split}
subspaces remains orthomodular, but ceases to be a lattice. This is
more specific than saying that some closed subspaces lack
orthogonal complements.
\end{remark}

At the level of inner-product geometry, $\Gamma=0$ gives the ordinary
Hilbert space and its complete closed-subspace lattice even in infinite
dimension. The earlier comparison of valuation and norm topologies
assumed $\Gamma\ne0$ and is not being applied to this degenerate case.
Thus nontrivial scale structure is an essential hypothesis of the
dichotomy, not just a notational preference.

\section{Fredholm perturbations reduce to a finite defect matrix}
\label{sec:schur}

\subsection{The hypotheses and the reduction}
Let $T_0:H\to G$ be an ordinary bounded Fredholm operator. This means
that its ordinary range is closed and that its kernel and cokernel
are finite dimensional. Write
\[
 N=\ker T_0,\quad X=N^\perp,\qquad
 Y=\ran T_0,\quad C=Y^\perp=\ker T_0^*.
\]
Let $m=\dim N$, $n=\dim C$. The restriction
$B_0=T_0|_X:X\to Y$ is a bounded isomorphism. We permit zero
summands; identity and inverse formulas on a zero space are
interpreted in their unique empty-block sense.

Suppose
\[
 A=T_0+E\in\BB{H}{G},\qquad E=0\text{ or }v(E)>0.
\]
Relative to $H=X\oplus N$, $G=Y\oplus C$, write
\begin{equation}\label{eq:blocks}
 A=\begin{pmatrix}B&C_1\\D&F_1\end{pmatrix}.
\end{equation}
Here $B=B_0+\text{positive-order terms}$ is invertible by the
Neumann formula, and the other three blocks have positive order
or are zero. Define the finite \emph{Schur defect matrix}
\begin{equation}\label{eq:schur}
 S=F_1-DB^{-1}C_1:\HH{N}\longrightarrow\HH{C}.
\end{equation}
Choosing ordinary orthonormal bases identifies it with an
$n\times m$ matrix over $\K$. Its entries have positive valuation
or are zero.

\begin{theorem}[Exact finite-defect reduction]\label{thm:schur}
With these hypotheses, the integral block operators
\begin{equation}\label{eq:LR}
 L=\begin{pmatrix}I&0\\-DB^{-1}&I\end{pmatrix},\qquad
 R=\begin{pmatrix}I&-B^{-1}C_1\\0&I\end{pmatrix}
\end{equation}
have integral inverses, are near the identity, and satisfy
\begin{equation}\label{eq:schuridentity}
 LAR=\begin{pmatrix}B&0\\0&S\end{pmatrix}=:\mathcal D.
\end{equation}
For $r=\rank_\K S$,
\begin{align}
 \dim_\K\ker A&=m-r,\label{eq:kerdim}\\
 \dim_\K(\HH{G}/\ran A)&=n-r,\label{eq:cokerdim}\\
 \ind_\K A&=m-n=\ind_\C T_0.\label{eq:index}
\end{align}
All identities are exact Hahn identities at arbitrary rank.
\end{theorem}
\begin{proof}
Factor $B=B_0(I+B_0^{-1}(B-B_0))$. Its inverse belongs to
$\BB{Y}{X}$, has integral valuation, and has residue $B_0^{-1}$.
The off-diagonal blocks in \eqref{eq:LR} are positive order.
Changing their signs gives the respective inverses of $L,R$.
Direct multiplication first clears the upper-right block by $R$
and then the lower-left block by $L$, leaving precisely
$F_1-DB^{-1}C_1$ in the lower-right corner. This proves
\eqref{eq:schuridentity}.

Because $L,R$ are bijections and $B$ is an isomorphism, the kernel
and cokernel of $A$ are isomorphic to those of the finite matrix $S$.
Finite rank-nullity gives \eqref{eq:kerdim}--\eqref{eq:index}.
For support validity, all inverse expansions use the positive
monoid generated by the support of $E$; multiplication and the
finite block decompositions preserve well-ordered support. No
analytic or norm-convergent infinite process is involved.
\end{proof}

This is a standard Schur-complement calculation applied in the
specified Hahn algebra. Its role is to isolate a finite part where
orthogonal linear algebra is safe. In particular, ``Fredholm'' in
\eqref{eq:index} is the finite algebraic kernel/cokernel condition
on these named spaces; it is not a compactoid-operator definition in
a different non-Archimedean category.

\subsection{Finite Moore--Penrose geometry without field closure}
A finite matrix $S:\K^m\to\K^n$ has an algebraic Moore--Penrose
inverse over $\K$ with its standard positive Hermitian form.
Indeed finite Gram--Schmidt gives
\[
 \K^m=\ker S\oplus(\ker S)^\perp,\qquad
 \K^n=\ran S\oplus(\ran S)^\perp.
\]
The restriction of $S$ between $(\ker S)^\perp$ and $\ran S$ is a
bijection of finite-dimensional spaces. Invert it there and set it
to zero on $(\ran S)^\perp$; call the resulting map $S^\dagger$.
It satisfies
\begin{equation}\label{eq:penrose}
 SS^\dagger S=S,\quad S^\dagger SS^\dagger=S^\dagger,
 \quad (SS^\dagger)^*=SS^\dagger,
 \quad (S^\dagger S)^*=S^\dagger S.
\end{equation}
This construction does not require that $\K$ be algebraically
closed or that $\F$ be real closed. In particular no assumption of
divisibility of $\Gamma$ has been inserted: the vector-norm roots
needed by Gram--Schmidt already exist by Proposition~\ref{prop:inner}.
These finite identities and their geometry are classical
Moore--Penrose linear algebra, not proposed new results.

\section{Exact Moore--Penrose inverses and least squares}
\label{sec:least}

\subsection{Finite defects restore orthogonal range decomposition}

\begin{theorem}[Fredholm-residue least-squares theorem]
\label{thm:MP}
For $A=T_0+E$ as in Section~\ref{sec:schur}, both $\ker A$ and
$\ran A$ are orthogonally split. There is a unique
$A^\dagger\in\BB{G}{H}$ satisfying the four Penrose identities
\eqref{eq:penrose} with $S$ replaced by $A$.
If $P$ projects onto $\ker A$ and $Q$ onto $\ran A$, it is explicitly
\begin{equation}\label{eq:MPformula}
 A^\dagger=(I-P)\,
 R\begin{pmatrix}B^{-1}&0\\0&S^\dagger\end{pmatrix}L\,Q.
\end{equation}
In particular
\begin{equation}\label{eq:MPproj}
 AA^\dagger=Q,\qquad A^\dagger A=I-P.
\end{equation}
\end{theorem}
\begin{proof}
The kernel is finite dimensional by Theorem~\ref{thm:schur}, so
its orthogonal projection $P$ exists. To prove the exact range
assertion rather than assume a Hilbert closed-range theorem, use
$A=L^{-1}\mathcal D R^{-1}$ from \eqref{eq:schuridentity}. Then
\[
 A^*=R^{-*}\mathcal D^*L^{-*},\qquad
 \ker A^*=L^*\ker\mathcal D^*.
\]
Here $L^{-*}=(L^*)^{-1}$ and similarly for $R$. On the finite Schur
block, ordinary finite orthogonal algebra over $\K$ gives
$(\ker S^*)^\perp=\ran S$. On the other block $B$ is invertible.
Therefore
\[
 (\ker\mathcal D^*)^\perp=\ran\mathcal D.
\]
For any subspace $Z$, the adjoint identity gives
$(L^*Z)^\perp=L^{-1}Z^\perp$. Consequently
\begin{equation}\label{eq:closedrangeexact}
 (\ker A^*)^\perp=L^{-1}\ran\mathcal D=\ran A.
\end{equation}
The subspace $\ker A^*$ has finite dimension $n-r$, so it is split.
Its orthogonal complement \eqref{eq:closedrangeexact} is split
as well; let $Q=I-P_{\ker A^*}$ be its projection. This proves
splitting using the finite reduction, not merely finite codimension.

Put
\[
 G_0=R\begin{pmatrix}B^{-1}&0\\0&S^\dagger\end{pmatrix}L.
\]
The block identity gives $AG_0A=A$, so $AG_0$ is the identity on
$\ran A$. Since $AP=0$, the map
$F=(I-P)G_0Q$ satisfies $AF=Q$. Moreover for every $x$,
$A(G_0Ax-x)=0$, so
\[
 FAx=(I-P)G_0Ax=(I-P)x.
\]
Thus $FA=I-P$. The range of $F$ lies in $(\ker A)^\perp$, and
$F=FQ$, which yields $FAF=F$. The two product identities give
self-adjointness of $AF,FA$ and hence all four Penrose equations.
Every factor in $F$ is adjointable: the finite Schur inverse and
finite-rank projections are Hahn operators, and all remaining
factors were already constructed in the Hahn algebra.

For uniqueness, any Penrose inverse $F'$ makes $AF'$ a
self-adjoint idempotent with range $\ran A$, hence $AF'=Q$.
Similarly $F'A$ is a self-adjoint idempotent with kernel $\ker A$,
hence $F'A=I-P$. Also $F'=F'AF'$, so its values lie in
$(\ker A)^\perp$. Both $F'f$ and $Ff$ lie there and solve
$Ax=Qf$; their difference lies in $\ker A$, and is therefore zero.
\end{proof}

\begin{corollary}[All least-squares and minimum-norm solutions]
\label{cor:LS}
For every $f\in\HH{G}$ let $u_0=A^\dagger f$. Then
\begin{equation}\label{eq:LSidentity}
 \norm{Ax-f}^2=\norm{(I-Q)f}^2+\norm{A(x-u_0)}^2
 \qquad(x\in\HH{H}).
\end{equation}
The least-squares minimizers are exactly $u_0+\ker A$.
Among them $u_0$ is the unique vector of minimum norm. The equation
$Ax=f$ is solvable exactly when $(I-Q)f=0$.
\end{corollary}
\begin{proof}
Write $Ax-f=A(x-u_0)-(I-Q)f$. Its two terms lie in the orthogonal
spaces $\ran A$ and $\ker A^*$, respectively. Pythagoras proves
\eqref{eq:LSidentity}; positivity identifies equality with
$A(x-u_0)=0$. Since $u_0\in(\ker A)^\perp$, one also has
$\norm{u_0+n}^2=\norm{u_0}^2+\norm{n}^2$ for $n\in\ker A$.
This proves the minimum-norm and exact-solvability assertions.
\end{proof}

\begin{remark}[Why this does not contradict the graph obstruction]
Example~\ref{ex:badgraph} has an infinite-dimensional ordinary range
defect, invisible to a finite matrix: its residue range is dense
and not closed. Theorem~\ref{thm:MP} starts with an ordinarily closed
Fredholm range and leaves only finitely many kernel and cokernel
coordinates to resolve. A finite matrix over $\K$ always has the
orthogonal decompositions needed here. The two results therefore
identify different regimes, not competing definitions of solvability.
\end{remark}

\section{The exact amplification scale from finite minors}
\label{sec:valuation}

\subsection{Finite valuation-ring reduction}
The valuation ring $\OO$ need not be Noetherian or discretely valued.
Nevertheless finite matrices admit the following elementary reduction.
We supply its proof to make the generality of the inverse-scale
formula explicit.

\begin{lemma}[Finite diagonal reduction over the valuation ring]
\label{lem:smith}
Let $S\in M_{n\times m}(\OO)$ have rank $r>0$. There are
$U\in\GL_n(\OO)$, $V\in\GL_m(\OO)$, both with inverses over $\OO$,
and exponents
$0\le s_1\le\cdots\le s_r$ such that
\begin{equation}\label{eq:smith}
 USV=\operatorname{diag}(t^{s_1},\ldots,t^{s_r},0).
\end{equation}
If every nonzero entry of $S$ has positive valuation, then $s_1>0$.
For $1\le j\le r$, set
\[
 \nu_j(S)=\min\{v(\det S_{I,J}):\abs I=\abs J=j,
                         \det S_{I,J}\ne0\},\qquad \nu_0(S)=0.
\]
Then
\begin{equation}\label{eq:minors}
 \nu_j(S)=s_1+\cdots+s_j.
\end{equation}
\end{lemma}
\begin{proof}
Choose a nonzero entry of minimum valuation and permute it to the
upper-left corner. Write it as $p=t^s u$ with $u\in\OO^\times$.
Multiplying its row by $u^{-1}$ makes the pivot $t^s$; all entries
still have valuation at least $s$. Because the pivot has minimum
valuation, it divides every entry in $\OO$. Subtract multiples of
the first row to clear the first column, then multiples of the
first column to clear the first row. These are invertible integral
row and column operations. The remaining lower-right block still
has all entries of valuation at least $s$: each new entry is a
difference of terms with that property. Repeat on this finite block.
The process terminates after $r$ nonzero pivots and produces
nondecreasing exponents. This proves \eqref{eq:smith} and the
strict positivity assertion.

Cauchy--Binet expresses every $j$-minor after an integral left or
right multiplication as a finite sum of original $j$-minors times
integral minors of the multiplier. Thus its minimum nonzero minor
valuation cannot decrease. Applying the integral inverse gives
the reverse inequality, so $\nu_j$ is invariant under both
operations. For the diagonal matrix in \eqref{eq:smith}, the least
nonzero $j$-minor valuation is $s_1+\cdots+s_j$. This proves
\eqref{eq:minors} and also shows that the exponents are invariants.
\end{proof}

The diagonal exponents in this lemma are valuation invariants, not
asserted eigenvalue or singular-value exponents of an arbitrary
infinite operator. No spectral theorem is being used.

\subsection{A sharp formula for the Moore--Penrose inverse}

\begin{theorem}[Exact finite-defect amplification]\label{thm:scale}
Let $A=T_0+E$ and $S$ be as in Theorem~\ref{thm:MP}. If
$r=\rank S>0$, then
\begin{equation}\label{eq:scale}
 v(A^\dagger)=-s_r
       =-\bigl(\nu_r(S)-\nu_{r-1}(S)\bigr)<0,
\end{equation}
where the $s_j$ are the exponents in Lemma~\ref{lem:smith}.
If $S=0$ and $T_0\ne0$, then
\begin{equation}\label{eq:nolift}
 v(A^\dagger)=0,
 \qquad [t^0]A^\dagger=T_0^\dagger.
\end{equation}
If $S=0=T_0$, then $A=A^\dagger=0$.
\end{theorem}
\begin{proof}
Combine the Schur reduction, multiplication of the regular block
by $B^{-1}$, and the integral finite transformations in
Lemma~\ref{lem:smith}. There are bounded-coefficient Hahn
isomorphisms $J$ on the codomain side and $K_1$ on the domain side,
with integral inverses, such that
\begin{equation}\label{eq:fullnormal}
 JAK_1=\mathcal E
  =\operatorname{diag}(I_X,t^{s_1},\ldots,t^{s_r},0).
\end{equation}
To write the regular identity $I_X$ we identify $Y$ with $X$ by
$B_0$ or $B$ as appropriate; this is an isomorphism between the
named coefficient summands, not a dimension-changing operation.
Every change and its inverse is integral. Zero regular summands
are allowed.

Assume first $r>0$. The diagonal map
\[
 Z_0=\operatorname{diag}(I_X,t^{-s_1},\ldots,t^{-s_r},0)
\]
is an inner generalized inverse of $\mathcal E$ and has valuation
$-s_r$. Therefore
$G=K_1Z_0J$ satisfies $AGA=A$ and has valuation $-s_r$, because
integral invertible multipliers preserve operator valuation.
The proof of Theorem~\ref{thm:MP} works with any such $G$ and gives
\[
 A^\dagger=(I-P)GQ.
\]
Both orthogonal projections are integral by
Lemma~\ref{lem:integralproj}. Hence
$v(A^\dagger)\ge-s_r$.

For the converse bound put
$Z=K_1^{-1}A^\dagger J^{-1}$. The identity
$AA^\dagger A=A$ implies $\mathcal E Z\mathcal E=\mathcal E$.
At the $j$th nonzero finite pivot, its diagonal entry says
\[
 t^{s_j}Z_{jj}t^{s_j}=t^{s_j},
 \qquad\text{so}\qquad Z_{jj}=t^{-s_j}.
\]
A nonzero scalar coordinate of valuation $-s_r$ forces
$v(Z)\le-s_r$. The integral invertible transformations preserve
valuation, so $v(Z)=v(A^\dagger)$. This proves equality; the
minor formula follows from \eqref{eq:minors}.

If $S=0$, the generalized inverse
$G_0=R\operatorname{diag}(B^{-1},0)L$ is integral and its residue
is $T_0^\dagger$ relative to the ordinary orthogonal decompositions.
Thus $A^\dagger=(I-P)G_0Q$ is integral. To identify the projection
residues, observe from the Schur formulas with $S=0$ that
\[
 \ker A=R(0\oplus\HH{N}),\qquad
 \ker A^*=L^*(0\oplus\HH{C}).
\]
Because $R,L^*$ and their inverses are near the identity, their
integral residue spaces are $N,C$. Corollary~\ref{cor:residue}
therefore gives $P_0=P_N$ and $Q_0=P_Y$. Taking residues in
\eqref{eq:MPformula} yields
$[t^0]A^\dagger=(I-P_N)T_0^\dagger P_Y=T_0^\dagger$.
This is nonzero if $T_0\ne0$, proving \eqref{eq:nolift}.
If $T_0=0$ is Fredholm, both $H,G$ are finite dimensional and the
regular block is empty; then $S=A$, so $S=0$ gives $A=0$.
\end{proof}

\begin{corollary}[Exactly when the inverse stays finite]
\label{cor:finiteinverse}
For a nonzero ordinary Fredholm $T_0$, the following are equivalent:
$A^\dagger$ is integral; $S=0$; $\dim_\K\ker A=\dim_\C\ker T_0$.
They are also equivalent to preservation of the ordinary cokernel
dimension. If even one finite defect direction is lifted, the
Moore--Penrose inverse acquires a strictly negative valuation,
exactly as measured by \eqref{eq:scale}.
\end{corollary}
\begin{proof}
Use \eqref{eq:kerdim}, \eqref{eq:cokerdim}, and
Theorem~\ref{thm:scale}; all $s_j$ are strictly positive.
\end{proof}

\subsection{An off-diagonal term can dominate the inverse scale}

\begin{example}[A fourth-order pole from second- and third-order diagonals]
\label{ex:pole}
Fix $\eta>0$, put $q=t^\eta$, and on $H_0\oplus\C^2$ take
\[
 T_0=I_{H_0}\oplus0_2,
 \qquad A=I_{H_0}\oplus
 S,\qquad S=\begin{pmatrix}q^2&q\\0&q^3\end{pmatrix}.
\]
Here $H_0$ may be infinite dimensional. The finite Schur matrix is
exactly $S$, with
\[
 \nu_1(S)=\eta,\qquad \nu_2(S)=5\eta,
 \qquad s_1=\eta,
 \quad s_2=4\eta.
\]
Since $A$ is bijective, $A^\dagger=A^{-1}$, and explicitly
\[
 S^{-1}=\begin{pmatrix}q^{-2}&-q^{-4}\\0&q^{-3}\end{pmatrix}.
\]
The inverse valuation is $-4\eta$, not $-2\eta$ or $-3\eta$.
Off-diagonal coupling is measured by minors and cannot be recovered
merely by inspecting the orders on the original diagonal.
\end{example}

More generally, for positive $\alpha,\beta,\gamma\in\Gamma$,
\[
 S=\begin{pmatrix}t^\alpha&t^\beta\\0&t^\gamma\end{pmatrix}
 \quad\Longrightarrow\quad
 v(S^{-1})=-\bigl(\alpha+\gamma-
                       \min\{\alpha,\beta,\gamma\}\bigr).
\]
This follows either from the two-minor formula or by writing its
inverse. The expression makes sense in any ordered abelian group.
For example, in $\Q\oplus\Q$ with lexicographic order take
$\alpha=(1,0)$, $\beta=(0,1)$, $\gamma=(1,1)$. Then the inverse
valuation is $-(2,0)$. There is no need to collapse these scales to
real numbers or choose an Archimedean embedding.

\section{Coherent extension of scales and the full surreal class}
\label{sec:global}

\subsection{Set-sized scalar extension preserves the formulas}
An ordered embedding $\Gamma\hookrightarrow\Delta$ sends a
well-ordered support to a well-ordered support, and defines embeddings
of scalar fields, vector spaces, and bounded-coefficient Hahn
operator algebras. It preserves products, conjugation, adjoints,
and the strong sums used in this article. A projection therefore
extends to a projection on the larger Hahn space.

\begin{proposition}[Coherence of the constructive results]
\label{prop:extension}
Under any ordered embedding $\Gamma\hookrightarrow\Delta$:
\begin{enumerate}[label=\textup{(\roman*)}]
\item the projection normal form and its direct rotation extend by
coefficient transport;
\item for a Fredholm-residue $A$, its transported Moore--Penrose
inverse is the Moore--Penrose inverse of the transported operator;
\item the Schur rank is unchanged, its nonzero minor valuations are
transported, and the exact amplification formula is unchanged under
that transport.
\end{enumerate}
No cofinality assumption on the embedding is required.
\end{proposition}
\begin{proof}
The normal form is made from finite algebraic expressions and
positive-support binomial series, all preserved by the embedding.
The transported Moore--Penrose inverse still satisfies the four
Penrose identities on the larger space, since they are operator
Hahn identities. Uniqueness in Theorem~\ref{thm:MP} identifies it.
A finite determinant is preserved by an injective field map, so the
largest nonzero minor size, namely the rank, is unchanged. Its
valuation is mapped by the ordered embedding, which preserves
minima of finite sets; this proves the final assertion.
\end{proof}

These are algebraic and support-coherence statements. They do not
assert that valuation completions or topological closures commute
with a noncofinal scale enlargement. The three-duals report explains
why those topological operations behave differently \cite{RepoDuals}.

\subsection{A literal vector space over all surcomplex scalars}
For readers who want all of $\No[i]$, rather than a specified set
subfield, there is a precise class formulation. Work here in
von Neumann--Bernays--G\"odel class theory without assuming global
choice, and keep the ordinary Hilbert spaces $H,G$ as sets. Define
\[
 \mathscr H_{\No}(H)=
 \left\{\sH_{\gamma\in S}h_\gamma t^\gamma:
 S\subset\No\text{ is a set and well ordered},\ h_\gamma\in H\right\},
 \qquad t^\gamma=\omega^{-\gamma}.
\]
Vectors can be encoded as sets of nonzero coefficient pairs. The
collection of all these vectors is a proper class when $H\ne0$.
By Conway normal forms the scalar class
$\C((t^{\No}))$ is exactly $\No[i]$ in this notation. The vector
operations and inner product remain the same coefficient formulas.
Any finite calculation involves a set of supports and lies in one
set-sized exponent subgroup. This supplies a literal positive
inner-product vector class over all surcomplex scalars, rather than
calling a proper class a set.

\begin{theorem}[Class extension and set support of adjoints]
\label{thm:class}
In this class formulation the following hold.
\begin{enumerate}[label=\textup{(\roman*)}]
\item The inner product, vector norm, and all support-admissible
set-indexed strong sums are well defined over $\No[i]$.
\item Every everywhere-defined surcomplex-linear map
$A:\mathscr H_{\No}(H)\to\mathscr H_{\No}(G)$ with an
everywhere-defined adjoint has a single \emph{set-supported}
Hahn expansion in ordinary bounded maps $H\to G$.
\item Every orthogonal projection, every amplified-graph criterion
of Theorem~\ref{thm:threshold}, and every Fredholm-residue
Moore--Penrose construction above applies on these full vector
classes. Each resulting operator has its coefficients in one
set-sized Hahn workspace, although it acts on vectors from all
workspaces.
\end{enumerate}
\end{theorem}
\begin{proof}
For (i), each operand has set support. The sum, product, and inner
product use only finite unions, support sums, and finite coefficient
convolutions of those sets. Their validity is the already proved
set-sized argument in the subgroup they generate. A set-indexed
admissible family has set union of supports by replacement and union;
the same reduction applies to its strong sum. The positive-root
construction for norms also has set support.

For (ii), restrict $A$ to the set $H$ of constant vectors. Class
replacement says that the image $\{Au:u\in H\}$ is a set. Therefore
\[
 S_A=\bigcup_{u\in H}\supp(Au)
\]
is a set of surreal exponents. Do the same for the adjoint restricted
to $G$. Extract coefficients as in Theorem~\ref{thm:adjoint}; the
ordinary closed graph argument makes them bounded, and the Baire
argument makes the nonzero coefficient support well ordered.
That proof only uses the two set coefficient spaces and the support
sets just exhibited. Finally, testing against constants identifies
the convolution expansion on every class vector. This proves one
set-supported expansion for the entire operator, not a separate
expansion for each input.

For (iii), an orthogonal projection is everywhere adjointable, so
(ii) applies and the integral-projection and direct-rotation proofs
use its set support. For a given surcomplex scalar $a$ choose a
set-sized group containing its normal-form support; the residue
proof of the graph threshold applies without alteration, as does
its explicit inverse in the closed-range case. For a positive-order
adjointable perturbation $E$, part (ii) places all its coefficients
in one set-sized group. The Schur and Moore--Penrose formulas can
be formed there. Their identities then hold on every class vector:
include that vector's support in a larger set-sized group and apply
Proposition~\ref{prop:extension}. Necessity in the graph obstruction
also persists, since any hypothetical global projection would have
an ordinarily closed residue by the same argument as before.
\end{proof}

This extension explains why the operator theorems genuinely concern
surcomplex scalars, not only a convenient notation for ordinary
series. It does not turn the global vector class into an ordinary
set-sized Hilbert space. Nor do we form a ``class of all subspace
classes.'' Orthogonal projections can instead be represented by
their set Hahn codes; the non-meet example of
Theorem~\ref{thm:lattice} still obstructs the corresponding
projection-code order in infinite coefficient dimension.

\subsection{Global fine completeness has a different meaning}

\begin{proposition}[Set-indexed fine Cauchy nets stabilize]
\label{prop:globalnets}
In $\mathscr H_{\No}(H)$, a net indexed by a set that is Cauchy for
\emph{all} positive surreal norm radii is eventually constant.
A convergent set-indexed net in the same fine topology is likewise
eventually constant.
\end{proposition}
\begin{proof}
The nonzero pairwise differences of a set-indexed net form a set of
vectors, and their valuations form a set of surreal exponents.
Every set of surreals has a surreal strict upper bound: for a set
$S$, the Conway cut $\{S\mid\varnothing\}$ is larger than every
member. Choose $\rho$ above all these difference valuations.
The Cauchy condition at a sufficiently small monomial radius makes
all differences in a tail have valuation greater than $\rho$.
No nonzero difference can do so, hence that tail is constant.
For a convergent net use instead the set of nonzero differences
from its limit and the same argument.
\end{proof}

This is consistent with nonconstant convergent sequences such as
$t^n\to0$ inside $\C((t^\Q))$: those sequences only pass the
\emph{workspace} tests. In the full surreal field, a radius smaller
than all their nonzero terms prevents convergence. Thus the
spherically complete set-valued model, the class-valued fine model,
and strong formal summation should not be conflated. We make no
claim about arbitrary class-indexed Cauchy families.

\section{Scope, provenance, and questions left by the theory}
\label{sec:provenance}

\subsection{Repository comparison at a fixed revision}
The comparison used revision
\texttt{9a385d3957bdfe3d9ea79f9a524751c90bd2c894}, retrieved on
September 23, 2026. In particular the following repository results
are predecessors, not contributions of this article.

The \emph{Infinite-Dimensional Hahn Spectral Theory} report
\cite{RepoSpectral} already constructs $H((t^\Gamma))$, proves its
positive inner product and the automatic adjoint theorem, and studies
spectra and range defects. Its relevant labels include
\texttt{ihs:hh:prop:inner}, \texttt{ihs:hh:thm:adjoint}, and
\texttt{ihs:hh:thm:defect}. Its later Fredholm-window part uses the
classical direct rotation for isolated finite spectral clusters and
develops a coefficientwise Fredholm determinant.
The \emph{Three Duals at Surreal Scales} report \cite{RepoDuals}
separates represented, strong, and continuous duals and distinguishes
algebraic scalar extension, valuation completion, and strong closure.
It also exhibits failures of Hilbert-style orthogonal geometry.

The present article does not reproduce those spectral classifications,
claim a new discovery of nonrepresented continuous functionals, or
rename a Fredholm determinant as a Moore--Penrose inverse. It uses
the intrinsic adjoint algebra to characterize all split subspaces,
then applies that characterization to graphs, the projection poset,
and finite-defect best approximation. The class argument in
Theorem~\ref{thm:class} adds the explicit set-support descent step
when using all surcomplex scalars.

The relevant repository portions were inspected, not the complete
proof corpus or all unintegrated source manuscripts. The repository
itself distinguishes drafted mathematics, proof review, and formal
verification. Accordingly no claim of novelty relative to every
uninspected source, or of inherited Lean verification, is made here.

\subsection{Primary-literature comparison and novelty boundary}
The literature search was targeted, not exhaustive. It included the
phrases ``Hahn Hilbert orthogonal graph,'' ``Hahn Moore--Penrose,''
``Hahn Fredholm least squares,'' non-Archimedean orthogonal
complementation, and the named works below. Failure of those searches
to identify an exact predecessor is evidence for a candidate
contribution, not a proof that none exists under other terminology.

Aguayo--Nova study non-Archimedean Hilbert-like $c_0$ spaces and
orthogonal-complementation obstructions \cite{AN}. Their model is
not the layered coefficient space retaining all of an ordinary
Hilbert $H$. The comparison used the publication record and
available author-posted text; a blocked publisher download was not
silently treated as a completed full-text audit.

Bursztyn--Waldmann develop formal deformations of Hermitian vector
bundles and finite projective modules, including uniqueness up to
isometric equivalence \cite{BW}. Simon discusses unitaries relating
projections and records the classical direct-rotation formula
\cite{Simon}. Thus neither projection deformation nor that formula
is claimed to originate here.

Avrachenkov--Lasserre investigate analytic perturbations of
generalized inverses \cite{Avrachenkov}. M\"uller--Strohmaier develop
Hahn-meromorphic functions and a holomorphic Fredholm theorem
\cite{MS}. They provide important analytic precedents. The present
proofs use instead unrestricted well-ordered supports over an
arbitrary ordered group, finite-defect orthogonal geometry, and
exact leading valuation from minors. Their analyticity and
convergence hypotheses are not imported, and neither their
methods nor classical generalized inverses are presented as new.
The comparison for \cite{Avrachenkov} used its abstract and the
coauthor's publication description, not an unavailable full PDF.

\begin{table}[htbp]
\centering\small
\begin{tabularx}{\textwidth}{@{}p{.27\textwidth}YY@{}}
\toprule
Result here & Prior input & Proposed contribution or status\\
\midrule
Coefficient-Hilbert construction and automatic adjoints
 & Repository spectral report
 & Reproved background; not new\\
Residue normal form
 & Automatic adjoints; classical direct rotation
 & Global classification of split subspaces in this model\\
Amplified graphs
 & Normal equations; ordinary closed-range inverse
 & Exact finite/infinite threshold and residue obstruction\\
Projection-code order
 & Orthomodular identities; finite projection theorem
 & Lattice if and only if coefficient dimension is finite\\
Fredholm least squares
 & Schur reduction; finite Penrose geometry
 & Arbitrary-rank existence and exact infinite-space solutions\\
Sharp inverse valuation
 & Finite valuation-ring elimination; determinantal ideals
 & Exact finite-Schur minor formula for the full pseudoinverse\\
Full surcomplex extension
 & Class replacement; normal forms; automatic adjoints
 & Explicit set support of global adjointable operators\\
\bottomrule
\end{tabularx}
\caption{Attribution and the boundary of the proposed theorem package.}
\end{table}

\subsection{What the theory deliberately leaves open}
The Fredholm assumption is sufficient, not asserted necessary, for
an adjointable Moore--Penrose inverse. A broader characterization
would have to treat infinitely many interacting defect directions
without assuming that their inverses share a well-ordered support
of bounded coefficient operators. The graph threshold exhibits a
precise obstruction that any such characterization must explain.

Likewise, closed residue alone does not classify arbitrary
valuation-closed subspaces. The graph chart characterizes the split
ones, but deciding when an independently presented closed subspace
admits such a chart is a separate question. General tensor products,
frame bounds, spectral measures, unbounded operators with proper
domains, and infinite-family projection joins are not developed
here. In particular no classical infinite-dimensional spectral
theorem, measurement postulate, or physical interpretation is
claimed merely from the existence of a positive surreal norm.

For computational use, the finite Schur matrix is exact but its
entries may themselves be infinite Hahn series. The minor formula
specifies an invariant and an algebraic reduction; it is not an
algorithm for deciding equality or leading terms of arbitrary
uncomputably presented surreal numbers. Effective implementations
need an explicit support and coefficient representation.

\subsection{Proof status and reproducibility}
All central arguments are given in full above. Their main dependencies
are the ordinary Hilbert closed-graph, Baire, and bounded-inverse
theorems; the Hahn--Neumann support lemma; Conway normal forms;
and finite linear algebra. The full-class section additionally
uses class replacement in the stated class theory. No proof
assistant was run, and the manuscript has not been independently
refereed. The mathematical claims should be evaluated by checking
the proofs and their hypotheses, not the novelty wording.

The accompanying \texttt{checks.py} uses exact symbolic algebra to
test representative finite identities: the projection rotation,
a Schur factorization, four Moore--Penrose equations, and the
coupled-pole example. It also checks finite triangular exponent
examples for the minor formula. These checks can catch algebraic
sign or transpose mistakes; they do not establish support
well-ordering, infinite-dimensional boundedness, spherical
completeness, the class argument, or bibliographic priority.
The archive includes the source, compiled PDF, check output, and
a short proof-audit ledger. It contains no repository modifications.

\clearpage
\appendix
\section{A dependency and hypothesis checklist}

\begin{longtable}{@{}p{.24\textwidth}p{.35\textwidth}p{.33\textwidth}@{}}
\toprule
Statement & Essential hypotheses & Proof mechanism\\
\midrule
\endfirsthead
\toprule
Statement & Essential hypotheses & Proof mechanism\\
\midrule
\endhead
Inner product and norms & Set support; ordinary positive Hilbert form
 & Leading coefficient and binomial root\\[.5em]
Spherical completeness & Set-sized group and nested set of balls
 & Compatible fixed truncations\\[.5em]
Automatic adjoints & Everywhere-defined adjoint; ordinary complete coefficient spaces
 & Coefficient closed graph, Baire, separation by constants\\[.5em]
Projection normal form & Self-adjoint idempotent; ordinary Hilbert coefficients
 & Integral leading projection and direct rotation\\[.5em]
Infinite graph threshold & Ordinary bounded $T$; scalar of negative valuation
 & Residue $\ker T\oplus\ran T$ and closed-range inverse\\[.5em]
Lattice dichotomy & Nontrivial group; ordinary Hilbert coefficient space
 & Three-coordinate graph intersection; all lines split\\[.5em]
Fredholm least squares & Ordinary closed range; finite kernel and cokernel; positive-order global perturbation
 & Schur reduction, finite orthogonal complements, Penrose compression\\[.5em]
Exact inverse scale & Nonzero finite Schur rank
 & Integral diagonal reduction and a forced inverse pivot\\[.5em]
Scalar extension & Ordered group embedding; no cofinality needed
 & Transport of finite algebra and admissible strong sums\\[.5em]
Full-class adjoints & Set coefficient spaces; class replacement; everywhere-defined adjoint
 & Set image of constants, then the automatic adjoint proof\\[.5em]
Fine-net stabilization & Set index and all surreal radii
 & A surreal upper bound on all difference valuations\\
\bottomrule
\end{longtable}

No theorem in the main set-sized treatment assumes a divisible group.
No infinite-dimensional theorem uses a finite matrix spectral theorem
to diagonalize an arbitrary Hahn operator. No unqualified exchange
of an ordinary Hilbert coordinate sum with a strong Hahn family is
performed. All nearest-point statements compare exact ordered-field
norms; they do not replace minimization by an unproved infimum.

\clearpage
\begin{thebibliography}{99}

\bibitem{Conway}
J.~H. Conway,
\emph{On Numbers and Games}, second edition,
A K Peters, 2001.
\href{https://www.routledge.com/On-Numbers-and-Games/Conway/p/book/9781568811277}{Publisher record, ISBN 9781568811277}.

\bibitem{RepoSpectral}
Surreal project, repository maintained by V.~Reshetnikov,
\emph{Infinite-Dimensional Hahn Spectral Theory},
AI-assisted research report, September 22, 2026;
revision \texttt{9a385d3957bdfe3d9ea79f9a524751c90bd2c894}.
\href{https://github.com/VladimirReshetnikov/Surreal/blob/9a385d3957bdfe3d9ea79f9a524751c90bd2c894/docs/surcomplex/infinite-dimensional-hahn-spectral-theory/article.tex}{Pinned manuscript source}.

\bibitem{RepoDuals}
Surreal project, repository maintained by V.~Reshetnikov,
\emph{Three Duals at Surreal Scales: Strong Hahn Operators,
Completion Defects, and Invisible Functionals},
AI-assisted research report, September 22, 2026; same pinned revision.
\href{https://github.com/VladimirReshetnikov/Surreal/blob/9a385d3957bdfe3d9ea79f9a524751c90bd2c894/docs/surcomplex/three-duals-of-hahn-vector-spaces/article.tex}{Pinned manuscript source}.

\bibitem{RepoFound}
Surreal project, repository maintained by V.~Reshetnikov,
\emph{Foundations} report and collection's notation and reading guide;
same pinned revision, accessed September 23, 2026.
\href{https://github.com/VladimirReshetnikov/Surreal/tree/9a385d3957bdfe3d9ea79f9a524751c90bd2c894/docs/foundations-and-computation/foundations}{Pinned foundations directory}.
The collection guide records the fine-topology and set-size distinctions.

\bibitem{Simon}
B.~Simon,
\emph{Unitaries Permuting Two Orthogonal Projections},
arXiv:1703.05437, 2017, version 2.
\url{https://arxiv.org/abs/1703.05437}.
In particular, formula (4) records the classical direct rotation.

\bibitem{BW}
H.~Bursztyn and S.~Waldmann,
\emph{Deformation Quantization of Hermitian Vector Bundles},
arXiv:math/0009170, 2000.
\url{https://arxiv.org/abs/math/0009170}.

\bibitem{AN}
J.~Aguayo and M.~Nova,
\emph{Non-Archimedean Hilbert like spaces},
Bulletin of the Belgian Mathematical Society--Simon Stevin
\textbf{14} (2007), no.~5, 787--797.
\href{https://doi.org/10.36045/bbms/1197908895}{doi:10.36045/bbms/1197908895}.

\bibitem{Avrachenkov}
K.~E. Avrachenkov and J.-B. Lasserre,
\emph{Analytic perturbation of generalized inverses},
Linear Algebra and its Applications \textbf{438} (2013), no.~4,
1793--1813. Author-deposited record:
\url{https://inria.hal.science/hal-00926609}.
Coauthor's description:
\url{https://www.laas.fr/fr/homepages/lasserre/softwares/singular-perturbations/}.

\bibitem{MS}
J.~M\"uller and A.~Strohmaier,
\emph{The theory of Hahn meromorphic functions, a holomorphic
Fredholm theorem and its applications},
Analysis \& PDE \textbf{7} (2014), no.~3, 745--770.
\href{https://doi.org/10.2140/apde.2014.7.745}{doi:10.2140/apde.2014.7.745};
\url{https://arxiv.org/abs/1205.0236}.

\end{thebibliography}
\end{document}
